\begin{proof}
\begin{enumerate}[label=\arabic*.]
\item The linear orders divide the affine multiplicities and are
therefore coprime by \ref{maps:src:F.2.7}.  The fixed-rank calculation
of \ref{maps:src:F.2.9}, which uses only \ref{maps:src:F.3.1} and the
evenness in \ref{maps:src:F.2.15}, gives fixed rank two at each finite
generator.  An identity generator would make the product finite and
unipotent, hence the identity, contrary to the rank-three cusp.

\item Suppose the cohomology generators are $A,B$ with $A^2=I$.
The order of $B$ is odd.  The $(-1)$-eigenspace of $A$ has dimension
four, whereas $\ker(AB-I)$ has dimension three.  Their intersection
in the six-dimensional space contains a nonzero vector $v$.
Then $ABv=v$ and $Av=-v$ imply
\[
 Bv=Av=-v,
\]
which contradicts the odd order of $B$.  Interchanging the two
generators proves the same assertion for either finite site.

\item Delete one and two from the complete finite-order list
\eqref{maps:eq:F.3} and enumerate its coprime unordered pairs.
They are exactly the eight pairs in \eqref{maps:eq:F.28}.
Only the divisibility of the linear orders into the affine
multiplicities has been used.
\end{enumerate}
\end{proof}
